\section{Schedule management}
\label{sec:schedule}

The campaign schedule turns the WBS into a dated baseline for board approval. Dates come from a finish-to-start activity-on-node network. The pad calendar is seven days a week from 10~August to 1~December~2026, because LOX load, static fire and lift-off do not stop for weekends. Approved lift-off ($T$-0) stays 15~November~2026. That day still sits inside the charter's last-allowed window of Q1~2027. A five-day pull-in to 10~November exists only if R-14 is realised (ASA flight authorisation or a CASA NOTAM refuses the 15~November slot). It is not the baseline.

\subsection{Network logic and float analysis}
\label{sec:schedule-network}

Links are finish-to-start with zero lag, drawn activity-on-node. Textbook discrete-time notation writes $EF = ES + D$ when $EF$ is exclusive. Every date on the Bowen calendar is inclusive. The identities used here are
\begin{equation}
ES_{i} = \max_{j \in P_{i}}(EF_{j}) + 1,
\qquad
EF_{i} = ES_{i} + D_{i} - 1,
\label{eq:fwd}
\end{equation}
\begin{equation}
LF_{i} = \min_{k \in S_{i}}(LS_{k}) - 1,
\qquad
LS_{i} = LF_{i} - D_{i} + 1,
\label{eq:bwd}
\end{equation}
$P_{i}$ and $S_{i}$ are the predecessor and successor sets of activity $i$. Where $P_{i}$ is empty, $ES$ is the campaign start. Total float and free float are
\begin{equation}
TF_{i} = LS_{i} - ES_{i} = LF_{i} - EF_{i},
\qquad
FF_{i} = \min_{k \in S_{i}}(ES_{k}) - EF_{i} - 1,
\label{eq:float}
\end{equation}
If $S_{i}$ is empty, $FF_{i} = TF_{i}$. Launch-critical means $TF_{i} = 0$ on a path that sets $T$-0 or closeout. Table~\ref{tab:precedence} is the full pass. It matches the locked register.

The forward pass gives two $TF = 0$ chains. They stay independent until hot fire. The vehicle path starts at the external \$0 freight wait (A-100; not a Gilmour launch-critical bar). Receiving (A-122) follows, then stacking (A-123), Wet Dress Rehearsal (A-132) and static fire (A-133). The permit path is ASA / CASA / GBRMPA permits (A-113) into static fire. RF calibration (A-125) is a finish-to-start predecessor of static fire, so the array is live before hot fire. It is not launch-critical. After static fire the tail from post-fire diagnostics (A-134) through closeout (A-145) is one critical chain. The PEP lock (A-111) finishes on 2~October by an imposed Must-Finish-On (Gate~1). It is not launch-critical. The Week~4 charter approved the document, not field mobilisation.

The paths meet in two staged merges. Wet Dress Rehearsal waits on stacking, cryogenic GSE (A-124), avionics power-on (A-131) and Juru monitors (A-114). That is the technical merge. Static fire then waits on the rehearsal, flight permits \emph{and} RF calibration. That is the regulatory and range-instrument merge. Launch Readiness Review (LRR, A-141) waits on post-fire diagnostics, not on static fire. Signing that flight-clearance gate on incomplete hot-fire data would break the hold-point logic of an AS9100D system \citep{sae2016}.

Feeder float is not interchangeable. Avionics has one day into the rehearsal. Cryogenic GSE has three. RF calibration has 39 days. It sits in August--September so it does not sit on the October pad peak. Moving it does not change 15~November. Cleanroom mobilisation (A-121) has $TF = 3$ but $FF = 0$, so delay is absorbed before $T$-0 only by consuming the GSE float. Total float protects $T$-0. Free float protects the successor. Figure~\ref{fig:network_logic} draws the links.

\begin{figure}[H]
\centering
\includegraphics[width=\linewidth]{network_logic.pdf}
\caption[Activity-on-node network]{Activity-on-node network for the Bowen campaign. Node bands are $ES$ / ID / $EF$, $LS$ / duration / $LF$, and $TF$. Crimson arrows are launch-critical finish-to-start links. The vehicle path and the permit path merge at static fire; RF calibration is the labelled south-routed feeder into the same merge; LRR waits on post-fire diagnostics, not on static fire. The PEP lock is Gate-1 constrained and is not launch-critical. Freight receipt is a grey dashed external \$0 driver, not a crimson Gilmour bar.}
\label{fig:network_logic}
\end{figure}

\begin{landscape}
\small
\setlength{\tabcolsep}{4.2pt}
\renewcommand{\arraystretch}{1.15}
\begin{longtable}{@{}l l l r l c c c c r r c@{}}
\caption[Activity precedence and float register]{Master campaign activity precedence and float register. Dates are inclusive. $FF = \min(ES_{\text{succ}}) - EF - 1$ days, or $FF = TF$ if there is no successor. The PEP lock $TF = 0$ is an imposed Gate~1 finish, not launch-critical. The freight wait $TF = 0$ is an external pin (CP $=$ Ext), not a Gilmour launch-critical bar.}
\label{tab:precedence}\\
\toprule
ID & WBS & Activity & $D$ & Pred. & ES & EF & LS & LF & TF & FF & CP \\
\midrule
\endfirsthead
\caption[]{Master campaign activity precedence and float register (continued).}\\
\toprule
ID & WBS & Activity & $D$ & Pred. & ES & EF & LS & LF & TF & FF & CP \\
\midrule
\endhead
\bottomrule
\endfoot
A-100 & EXT & Vehicle receipt (external driver) & 35 & --- & 10 Aug & 13 Sep & 10 Aug & 13 Sep & 0 & 0 & Ext \\
A-111 & 1.1.1 & PEP and integration baseline & 54 & --- & 10 Aug & 2 Oct & 10 Aug & 2 Oct & 0* & 0 & Gate \\
A-112 & 1.1.2 & EVM control stand-up & 7 & A-111 & 3 Oct & 9 Oct & 25 Nov & 1 Dec & 53 & 53 & N \\
A-113 & 1.1.3 & ASA / CASA / GBRMPA permits & 76 & --- & 10 Aug & 24 Oct & 10 Aug & 24 Oct & 0 & 0 & Y \\
A-114 & 1.1.4 & Juru monitor mobilisation & 62 & --- & 10 Aug & 10 Oct & 15 Aug & 15 Oct & 5 & 5 & N \\
A-121 & 1.2.1a & Site mobilisation and cleanroom & 19 & --- & 10 Aug & 28 Aug & 13 Aug & 31 Aug & 3 & 0 & N \\
A-122 & 1.2.1b & Receiving and unpacking & 8 & A-100, A-121 & 14 Sep & 21 Sep & 14 Sep & 21 Sep & 0 & 0 & Y \\
A-123 & 1.2.2 & Crane and vertical stacking & 24 & A-122 & 22 Sep & 15 Oct & 22 Sep & 15 Oct & 0 & 0 & Y \\
A-124 & 1.2.3 & Cryo GSE piping and checkouts & 45 & A-121 & 29 Aug & 12 Oct & 1 Sep & 15 Oct & 3 & 3 & N \\
A-125 & 1.2.4 & RF telemetry calibration & 18 & A-121 & 29 Aug & 15 Sep & 7 Oct & 24 Oct & 39 & 39 & N \\
A-131 & 1.3.1 & Avionics power-on sweeps & 23 & A-122 & 22 Sep & 14 Oct & 23 Sep & 15 Oct & 1 & 1 & N \\
A-132 & 1.3.2 & Wet Dress Rehearsal (cold) & 9 & A-114, A-123, A-124, A-131 & 16 Oct & 24 Oct & 16 Oct & 24 Oct & 0 & 0 & Y \\
A-133 & 1.3.3 & 5-second static fire & 6 & A-132, A-113, A-125 & 25 Oct & 30 Oct & 25 Oct & 30 Oct & 0 & 0 & Y \\
A-134 & 1.3.4 & Post-fire / pre-kitted spare & 6 & A-133 & 31 Oct & 5 Nov & 31 Oct & 5 Nov & 0 & 0 & Y \\
A-141 & 1.4.1 & LRR and RSO clearance & 5 & A-134 & 6 Nov & 10 Nov & 6 Nov & 10 Nov & 0 & 0 & Y \\
A-142 & 1.4.2 & Range exclusion and countdown & 4 & A-141 & 11 Nov & 14 Nov & 11 Nov & 14 Nov & 0 & 0 & Y \\
A-143 & 1.4.3 & Launch execution ($T$-0) & 1 & A-142 & 15 Nov & 15 Nov & 15 Nov & 15 Nov & 0 & 0 & Y \\
A-144 & 1.4.4 & Telemetry decrypt and handover & 5 & A-143 & 16 Nov & 20 Nov & 16 Nov & 20 Nov & 0 & 0 & Y \\
A-145 & 1.4.5 & Demobilisation and closeout & 11 & A-144 & 21 Nov & 1 Dec & 21 Nov & 1 Dec & 0 & 0 & Y \\
\end{longtable}
\end{landscape}

\subsection{Schedule baseline, Gantt visualisation and milestone integrity}
\label{sec:schedule-baseline}

The dates in Table~\ref{tab:precedence} \emph{are} the schedule baseline. Figure~\ref{fig:gantt_chart} plots the same $ES$/$EF$ values. Launch-critical bars are crimson. Remaining total float is a teal whisker to $LF$. Permitting, heritage, mobilisation and GSE run through August and September. October is the physical peak: stacking 15~October, WDR 16--24~October, static fire 25--30~October. November is the flight-clearance tail.

Two Week~4 collisions are closed here. The charter dated pad integration 15~October in Objective~1 and stacking as M-4 on 28~October. The PEP locks M-4 to stacking on \textbf{15~October~2026}. That restores the cold-flow window before hot fire. Handover ``within five days of lift-off'' was also dated M-8 at 22~November, seven days after a 15~November $T$-0. M-8 is now telemetry handover on \textbf{20~November}. The facility licence is already held, so permits gate static fire, not stacking. Treating flight permits as a stacking predecessor was an initiation error.

Three gates sit on the baseline. Gate~1 is the PEP lock on 2~October. Gate~2 is permits on 24~October. Without those instruments static fire cannot start. That path is already $TF = 0$. Gate~3 is LRR on 10~November. Crashing stacking without crashing the permit path therefore cannot pull $T$-0.

\begin{figure}[!htbp]
\centering
\includegraphics[width=\linewidth]{gantt_chart.pdf}
\caption[Baseline Gantt chart]{Baseline Gantt from 10~August to 1~December~2026, plotted from the same $ES$/$EF$ values as Table~\ref{tab:precedence}. Crimson bars are launch-critical. The hatched grey bar is the external \$0 freight wait. Teal whiskers show total float to $LF$. Vertical markers are M-2 (navy, Gate~1) and M-4 through M-9. Option~C is not drawn on this baseline.}
\label{fig:gantt_chart}
\end{figure}

\subsection{Schedule compression: crashing and fast-tracking}
\label{sec:schedule-compression}

Fast-tracking overlaps sequential work, usually at zero direct cost. Crashing adds resource on critical work and is paid as a cost slope. The Board question is whether $T$-0 can move from 15~November to 10~November if R-14 is realised (ASA / CASA refuses the 15~November slot). Five days must come off a $TF = 0$ path that actually sets lift-off. After the static-fire merge that path is the tail from diagnostics through lift-off. Crashing stacking by five days saves \textbf{zero} days on $T$-0 while permits still finish on 24~October. Overlapping stacking with cryogenic GSE also saves zero days. GSE already runs in parallel with stacking.

The crash-cost slope on a crashed activity is
\begin{equation}
\text{Cost slope} = \frac{\Delta C}{\Delta T}
= \frac{\text{crash cost} - \text{normal cost}}{\text{normal duration} - \text{crash duration}}.
\label{eq:slope}
\end{equation}
Table~\ref{tab:compression} and Figure~\ref{fig:compression_curve} compare three options that each claim five days. Only Option~C is both CPM-valid and acceptable against the LRR hold.

\begin{table}[H]
\centering
\small
\setlength{\tabcolsep}{3.5pt}
\caption[Schedule compression evaluation]{Schedule compression evaluation. All three options save five days on $T$-0 if they work as drawn. Baseline $T$-0 stays 15~November. Option~C is a draw on the \$78{,}890 reserve if R-14 fires, not an addition to the \$1{,}500{,}000 base.}
\label{tab:compression}
\begin{tabularx}{\textwidth}{@{}l >{\raggedright\arraybackslash}X c r r >{\raggedright\arraybackslash}X l@{}}
\toprule
Option & Activities modified & Days & Added cost & Slope & Secondary risks & Verdict \\
\midrule
A Pure fast-track & Diagnostics $\to$ LRR SS+1 (review overlaps diagnostics) & 5 & \$0 & \$0/d & Review before diagnostic close-out; LRU rework; AS9100D hold-point failure & Rejected \\
B Pure crash & Diagnostics crashed 6~$\to$~1~d & 5 & \$40{,}000 & \$8{,}000/d & Fatigue; skipped thermal/electrical screens; recycle into LRR & Rejected \\
C Hybrid (if R-14 fires) & Diagnostics 6~$\to$~1~d extra labour (\$21{,}000); LRR kept at 5~d; countdown kept at 4~d & 5 & \textbf{\$21{,}000} & \textbf{\$4{,}200/d} & Screens kept; R-15 4~d marine hold intact & Contingent \\
\bottomrule
\end{tabularx}
\end{table}

Decision: baseline $T$-0 stays 15~November. If R-14 fires, the Board draws Option~C and lift-off moves to 10~November. That is the only option that is valid on the network and acceptable on quality.

Option~A is a pure fast-track. It buys five days at \$0 by starting the review one day after diagnostics start. The review board would sit while those diagnostics were still open. AS9100D is the quality rule this campaign uses for hold points \citep{sae2016}. The LRR hold (HP-4) is sign-off of that review. A hold point is a stop: the next job cannot start until the planned checks are done and the hold is released. Option~A would let flight clearance start before post-fire data are complete. Signing that gate on incomplete data is a hold-point failure. Waiting to sign until diagnostics finish does not save five days; it creates rework. Fast-tracking the review is rejected.

Option~B is a pure crash. It compresses diagnostics from six days to one. The slope is $\$40{,}000/5 = \$8{,}000$ per day. A 24-hour diagnostic war room on a vehicle that has just fired is the work most likely to skip thermal and electrical screens. Crashing diagnostics without extra screens is rejected.

Option~C is the hybrid. It takes the five days only from diagnostics. Extra labour keeps the thermal and electrical screens that Option~B would skip: two specialists and two technicians for five days ($2\times\$165 + 2\times\$95 = \$520$/h $\times$ 8~h $=$ \$4{,}160/d, published as the same \$4{,}200/d unit used on the crawler) for \$21{,}000. The slope is $\$21{,}000/5 = \$4{,}200$ per day. Launch Readiness Review stays five days. Countdown stays four days, so the marine hold on R-15 is not consumed. If invoked, diagnostics sit on 31~October only, the review on 1--5~November, countdown on 6--9~November, and lift-off on \textbf{10~November}. The \$21{,}000 is drawn from the \$78{,}890 contingency only if R-14 is realised. It is not loaded into the \$1{,}500{,}000 base. Figure~\ref{fig:gantt_chart} is not redrawn. Feeder floats do not tighten: avionics stays at $TF = 1$, cryogenic GSE at $TF = 3$.

Option~C is chosen because it is cheaper than Option~B (\$4{,}200/d not \$8{,}000/d), it keeps the LRR hold that Option~A would break, and crashing stacking cannot move $T$-0 while permits still finish on 24~October.

\begin{figure}[H]
\centering
\includegraphics[width=0.98\linewidth]{compression_tradeoff.pdf}
\caption[Compression cost and risk trade-off]{Compression trade-off for a five-day pull-in of $T$-0. Left: direct added cost. Right: residual schedule and quality risk after treatment. Option~C is the only CPM-valid pull-in that keeps LRR at five days, countdown at four days, and pays \$4{,}200 per day rather than \$8{,}000.}
\label{fig:compression_curve}
\end{figure}
